% Separate geometric input; no restricted-Petri nonvanishing is claimed.
\subsection{Smooth determinant curves in a sufficiently positive range}

\begin{lemma}[First jets of the two finite pushforwards]
\label{lem:DE-A3-positive-pushforward-jets}
\status{Candidate geometric input; complete elementary proof written.}
Keep the mixed-surface, $\alpha_2$-torsor and line-bundle notation
of Lemma~\ref{lem:DE-A3-positive-two-bundles}. Assume that each of
$J$ and $JQ$ is the third power of a very ample line bundle on $C$.
Then $E_0=\nu_*J$ and $E_1=\nu_*(JQ)$ separate first jets at every
point of $T$:
\[
 H^0(T,E_i)\longrightarrow
 E_i\otimes_{\mathcal O_T}\mathcal O_T/\mathfrak m_t^2
 \quad\hbox{is surjective for every }t\in T(k).
                                                        \tag{DE-SJ1}
\]
In particular the hypothesis holds if $J=J_0^{3d}$, with $J_0$
ample and $d\ge3$. In the original polarization notation this
applies whenever $M$ itself descends through $\rho$, with descended
bundle $J_0$, and the section bundle is $M^{3d}$ with $d\ge3$.
\end{lemma}
\begin{proof}
The finite flat isogeny $\nu:C\to T$ is a torsor under the smooth
surface's nonzero invariant $\alpha_2$ direction. At its unique
geometric point over $t$, suitable completed regular parameters
identify the map with
\[
                  k[[x,z^2]]\longrightarrow k[[x,z]].
\]
This follows directly from the local torsor algebra
$\mathcal O_T[z]/(z^2-h)$: smoothness of $C$ forces $dh\ne0$
at the point, so $h$ is a regular parameter on $T$.
The inverse image of the first infinitesimal neighbourhood of $t$
therefore has ideal
\[
                   (x^2,xz^2,z^4)\subset k[[x,z]].
\]
Its quotient has basis $1,z,z^2,z^3,x,xz$, and this ideal contains
$(x,z)^4$. Thus it suffices that $J$ and $JQ$ separate third jets
at the corresponding point of $C$.

If $V$ is very ample, choose a section $v_0$ nonzero at the point
and sections whose ratios with $v_0$ give a basis of
$\mathfrak m/\mathfrak m^2$. Products of three sections of $V$,
using $v_0$ for the remaining factors, give every monomial of total
degree at most three modulo $\mathfrak m^4$. These monomials span
that local Artin algebra, independently of characteristic.
Consequently $V^3$ separates third jets. The completed calculation
computes the same finite local quotient before completion, proving
(DE-SJ1) by finite pushforward.

For the stated sufficient condition, $J_0^d$ is very ample for
$d\ge3$. Choose a cube root of $Q$ in $\Pic^0(C)(k)$; every
$\Pic^0$ twist of $J_0^d$ is a translate of it, since its
polarization isogeny is surjective on algebraically closed points.
Thus $JQ$ is also a cube of a very ample line bundle. Descent
preserves ampleness under the finite surjection $\rho$ and gives
$J=J_0^{3d}$ in the final situation.
\end{proof}

\begin{proposition}[A smooth determinant curve and a separable pencil]
\label{prop:DE-A3-positive-smooth-separable-pencil}
\status{Candidate geometric input; complete proof written.}
Under the preceding hypotheses, a general pair $(s_0,s_1)$ has
reduced disjoint zero schemes and a smooth integral determinant
curve $D=(s_0\wedge s_1=0)$. The logarithmic differential
$\eta_P|_D$ is nonzero. Hence the base-point-free pencil
$\langle a_0,a_1\rangle$ on $A=R_0^2$ is separable.

If a sharp-source open is nonempty for these bundles, the general
pair may also be chosen in that open. At such a pair the ordinary
curve multiplication map
\[
 H^0(D,A)\otimes H^0(D,\omega_D A^{-1})
       \longrightarrow H^0(D,\omega_D)
                                                        \tag{DE-SJ2}
\]
has rank at least three, and has rank four unless
$\omega_D A^{-2}$ is trivial.
\end{proposition}
\begin{proof}
First-jet separation implies that the value and its first
derivatives of each vector-bundle section can be chosen
independently at every point. The incidence where a rank-two
section vanishes and its derivative has rank less than two has
codimension at least three in the section space, over a surface
of dimension two. A general zero scheme is therefore finite and
reduced. Simultaneous vanishing of both sections imposes four
independent value conditions, so the two zero schemes are disjoint.

On this disjoint-zero open, a point of the determinant divisor has
matrix rank exactly one. Normalize one nonzero section value to
the first coordinate vector. The determinant and its two first
derivatives are then three independent conditions on the first
jet of the other section. Their incidence has codimension three
over the two-dimensional base. Thus a general determinant divisor
is smooth. Its line bundle $LP$ is ample, so it is nonempty and
connected: the restriction sequence and ample vanishing give
$H^0(D,\mathcal O_D)=k$. A smooth connected curve over the
algebraically closed field is integral.

Suppose $\eta_P$ restricted to $\Omega_D^1$ were zero. The conormal
sequence would lift its restriction as a section of
$\mathcal O_D(-D)$. This line bundle has negative degree, so has
no nonzero section. But the nonzero invariant differential
$\eta_P$ is nowhere zero as a section of $\Omega_T^1$, and hence
its restriction to $\Omega_T^1|_D$ is nonzero, a contradiction.
The already proved identity
$d\log(a_1/a_0)=\eta_P|_D$ now proves that the pencil is separable.

All conditions just established are nonempty opens in the
irreducible space of section pairs; they meet any nonempty
sharp-source open. There $h^0(A)=h^0(\omega_D A^{-1})=2$ by the
proved contraction theorem, and the two pencil sections form a
basis of $H^0(A)$. The base-point-free pencil sequence, tensored
by $\omega_D A^{-1}$, identifies the kernel of (DE-SJ2) with
\[
                          H^0(D,\omega_D A^{-2}).
\]
This line bundle has degree zero on the smooth integral curve.
Its section space has dimension zero unless it is trivial, in
which case the dimension is one. The domain of (DE-SJ2) has
dimension four, proving both assertions.
\end{proof}

\begin{remark}[The distinction from the restricted variation]
\label{rem:DE-A3-positive-smooth-pencil-scope}
\status{Candidate scope statement; no new sharp-rank conclusion.}
The preceding multiplication map is a curve-pencil map. No
identification of it with the actual restricted surface cup-wedge
map in (DE-PT3)--(DE-PT7) has been proved. Thus its rank bound does
not establish the required common-radical vanishing. The power
hypothesis is explicit: positivity of the ordinary index of
$M^d$ alone does not imply that $M$ itself descends along $\rho$.
In particular no unproved low-degree jet separation or determinant
curve smoothness is used in earlier arguments.
\end{remark}
